% Generated draft: requires mathematical and source review.
\documentclass[11pt]{article}
\usepackage[margin=25mm]{geometry}
\usepackage{fontspec}
\setmainfont{Latin Modern Roman}
\usepackage{amsmath,amssymb,mathtools,unicode-math}
\setmathfont{Latin Modern Math}
\usepackage{enumitem,needspace,xurl,hyperref}
\setlength{\parindent}{0pt}
\setlength{\parskip}{5pt}
\setlength{\emergencystretch}{3em}
\newcommand{\N}{\mathbb{N}}
\newcommand{\Z}{\mathbb{Z}}
\newcommand{\Q}{\mathbb{Q}}
\newcommand{\R}{\mathbb{R}}
\newcommand{\C}{\mathbb{C}}
\newcommand{\F}{\mathbb{F}}
\newcommand{\A}{\mathbb{A}}
\newcommand{\PP}{\mathbb{P}}
\newcommand{\E}{\mathbb{E}}
\newcommand{\eps}{\varepsilon}
\newcommand{\dd}{\,\mathrm{d}}
\begin{document}
\section*{UnsolvedMath ID: 10400064}
\textbf{Source problem code:} catalogue entry
\subsection{Definitions and mathematical statement}
A rational Drinfel'd associator is, in a standard normalization, a group-like formal series $\Phi(X,Y)\in\Q\langle\langle X,Y\rangle\rangle$, with constant term one, satisfying the associator pentagon and hexagon identities for a fixed rational associator parameter. Group-like refers to the completed Hopf algebra with primitive generators $X,Y$.

For example, the pentagon is expressed in the completed infinitesimal pure-braid algebra on four strands as
\[
\Phi(t_{12},t_{23}+t_{24})\Phi(t_{13}+t_{23},t_{34})
=\Phi(t_{23},t_{34})\Phi(t_{12}+t_{13},t_{24}+t_{34})\Phi(t_{12},t_{23}),
\]
where the symmetric generators $t_{ij}$ satisfy $[t_{ij},t_{kl}]=0$ for disjoint pairs and $[t_{ij},t_{ik}+t_{jk}]=0$ for distinct indices. The hexagon and parameter normalization must match the chosen convention; the pentagon alone is not supplied as the whole definition.

The problem attributed to J. Roberts is to construct such a rational associator \emph{in the context of rational homotopy theory}. The requested contribution is a construction in that setting with its identity/compatibility proofs, not merely an unspecified existence assertion. All coefficients must be rational and compatible in every degree; separate finite truncations or reduction of transcendental coefficients do not automatically provide it.

The source does not choose a rational-homotopy model, spaces or operadic structures, construction algorithm, parameter normalization, or criterion for what constitutes the intended contextual construction. No particular formality/minimal-model theorem or canonical associator is inferred from the brief request.
\subsection{Short English statement}
Construct an all-degree rational associator using a specified rational-homotopy framework, proving the required associator identities and normalization. The source does not identify that framework or construction criterion.
\subsection{Sources}
\begin{itemize}
\item Local UnsolvedMath registry, permanent ID 10400064, exact source code catalogue entry, Problem 3.14; source attribution J. Roberts, no proposer/year in metadata. Registry status partially solved, classification PARTIAL-PROGRESS; the registry title is truncated.
\item Ohtsuki, \emph{Problems on invariants of knots and 3-manifolds} (2002), Problem 3.14, PDF page 63, \url{https://msp.org/gtm/2002/04/gtm-2002-04-024s.pdf}. Supplied pdf-font-structure extraction accessed July 29, 2026; provenance distinguishes 231 bold numbered items. No PDF was accessed here.
\item The supplied summary mentions unspecified rational-homotopy associator constructions by different authors without names, titles or theorems, so no precise later resolution of this construction request is asserted.
\item Provenance flags current primary-literature status and redistribution terms NEEDS\_REVIEW.
\end{itemize}
\end{document}
